% =====================================================================
%  Раздел «Архитектура подсистемы контроля целостности».
%  Вставляется в architecture_description.tex (например, после раздела
%  «Подсистема контроля целостности») командой
%      % =====================================================================
%  Раздел «Архитектура подсистемы контроля целостности».
%  Вставляется в architecture_description.tex (например, после раздела
%  «Подсистема контроля целостности») командой
%      % =====================================================================
%  Раздел «Архитектура подсистемы контроля целостности».
%  Вставляется в architecture_description.tex (например, после раздела
%  «Подсистема контроля целостности») командой
%      % =====================================================================
%  Раздел «Архитектура подсистемы контроля целостности».
%  Вставляется в architecture_description.tex (например, после раздела
%  «Подсистема контроля целостности») командой
%      \input{integrity_architecture}
%  Рисунок: в draw.io открыть integrity_architecture.drawio,
%  «Файл -> Экспорт -> PDF» с галочкой «Обрезать», сохранить как
%  integrity_architecture.pdf рядом с документом.
% =====================================================================
\section{Архитектура подсистемы контроля целостности}\label{sec:integrity-arch}

Подсистема реализована процессом \texttt{fw-integrity} на Python~3.13,
используется только стандартная библиотека. Процесс работает от
\texttt{root}: ему нужно читать все контролируемые файлы, восстанавливать
их и останавливать службы в аварийном режиме. Схема приведена на
рисунке~\ref{fig:integrity-arch}, назначение компонентов~— в
таблице~\ref{tab:integrity-parts}.

\begin{landscape}
\begin{figure}[p]
\centering
\IfFileExists{integrity_architecture.pdf}{%
  \includegraphics[width=\linewidth,height=0.86\textheight,keepaspectratio]{integrity_architecture.pdf}}{%
\IfFileExists{integrity_architecture.png}{%
  \includegraphics[width=\linewidth,height=0.86\textheight,keepaspectratio]{integrity_architecture.png}}{%
  \fbox{\parbox{0.8\linewidth}{\centering\vspace{2cm}
  Экспортируйте схему из draw.io в файл \texttt{integrity\_architecture.pdf}.\vspace{2cm}}}}}
\caption{Архитектура подсистемы контроля целостности}\label{fig:integrity-arch}
\end{figure}
\end{landscape}

{\small
\begin{longtable}{|L{3.4cm}|L{2.6cm}|L{8.2cm}|}
\caption{Компоненты подсистемы контроля целостности}\label{tab:integrity-parts}\\
\hline
\textbf{Компонент} & \textbf{Модуль} & \textbf{Назначение}\\
\hline\endfirsthead
\hline
\textbf{Компонент} & \textbf{Модуль} & \textbf{Назначение}\\
\hline\endhead
API-сервер & \texttt{protocol} &
  UNIX-сокет \texttt{/run/fw/integrity.sock}, JSON-строки (АПИ-03, АПИ-04),
  коды ошибок (АПИ-09). Права проверяются по контексту вызова: ТЧ, ЧЗ или
  подсистема; отказ регистрируется событием~106\\ \hline
Планировщик & \texttt{daemon} &
  Проверка при запуске до сигнала готовности systemd, каждые
  1\,--\,1440~мин (по умолчанию 10) и по запросу (ЦЛС-03)\\ \hline
Ядро & \texttt{service} &
  Команды API, режимы «обычный» и «аварийный», легитимные изменения
  файлов (ЦЛС-02), период проверки. Проверка, восстановление и обновление
  эталонов выполняются по очереди под одной блокировкой\\ \hline
Проверка файлов & \texttt{checker}, \texttt{manifest} &
  SHA-256 каждого файла с эталоном; поиск изменённых, удалённых,
  посторонних (в том числе скрытых) и нечитаемых файлов (ЦЛС-01)\\ \hline
Проверка ОС & \texttt{os\_state} &
  Сверка с таблицей ОБР-02: пересылка ядром, IPv6, адреса портов, promisc
  и MTU, мосты и VLAN, аппаратные функции, правило \texttt{drop} на
  \texttt{netdev ingress}\\ \hline
Восстановление & \texttt{store} &
  Изменённый или удалённый файл восстанавливается из резервной копии
  после её проверки по эталону; посторонний файл перемещается в карантин
  (ЦЛС-05)\\ \hline
Глобальная блокировка & \texttt{actions} &
  Команда \texttt{set\_emergency} подсистеме фильтрации и остановка
  служб fw-filter, fw-web, ssh (ЦЛС-04); при выходе~— обратные действия
  (ЦЛС-06)\\ \hline
Отправка событий & \texttt{events} &
  Команда \texttt{log\_event} подсистеме регистрации (коды 416,
  500\,--\,505); очередь до 1000 событий, если регистрация недоступна\\ \hline
\end{longtable}
}

\textbf{Режимы работы.} В обычном режиме трафик обрабатывается, а файлы
проверяются по расписанию. Любое нарушение (или невозможность проверки:
нет эталонов, повреждён манифест) переводит подсистему в аварийный режим.
Сначала режим записывается в \texttt{state.json}, затем включается
глобальная блокировка, затем регистрируются события~501 и~502.
Возврат в обычный режим происходит только после повторной проверки без
нарушений (событие~504). Запуск fw-filter и fw-web защищён проверкой
\texttt{fw-integrity gate} в \texttt{ExecStartPre}: служба стартует, только
если в текущей загрузке уже была успешная проверка (ЦЛС-07).

\textbf{Легитимные изменения.} Подсистемы, которые пишут конфигурационные
файлы (политику, сохранённую конфигурацию, учётные записи), вызывают
\texttt{begin\_update}, атомарно записывают файл и вызывают
\texttt{commit\_update}. Пока изменение не завершено, файл не проверяется.
После \texttt{commit\_update} эталон и резервная копия обновляются
(событие~505). Изменение без этих вызовов считается нарушением.

\textbf{Данные подсистемы} хранятся в каталоге
\texttt{/var/lib/fw/integrity}, доступном только \texttt{root}:
\begin{itemize}
\item \texttt{reference.json}~— эталонные SHA-256, вид, размер и права файлов;
\item \texttt{backup/}~— резервные копии, имя файла~— его SHA-256;
\item \texttt{quarantine/}~— посторонние файлы, убранные при восстановлении;
\item \texttt{state.json}~— режим, перечень нарушений, период и
  идентификатор загрузки.
\end{itemize}
Перечень контролируемых файлов задаёт манифест \texttt{manifest.json}:
\texttt{static} для кода и неизменяемых файлов, \texttt{config} для
изменяемых конфигураций. Сам манифест и файл настроек всегда
контролируются.

%  Рисунок: в draw.io открыть integrity_architecture.drawio,
%  «Файл -> Экспорт -> PDF» с галочкой «Обрезать», сохранить как
%  integrity_architecture.pdf рядом с документом.
% =====================================================================
\section{Архитектура подсистемы контроля целостности}\label{sec:integrity-arch}

Подсистема реализована процессом \texttt{fw-integrity} на Python~3.13,
используется только стандартная библиотека. Процесс работает от
\texttt{root}: ему нужно читать все контролируемые файлы, восстанавливать
их и останавливать службы в аварийном режиме. Схема приведена на
рисунке~\ref{fig:integrity-arch}, назначение компонентов~— в
таблице~\ref{tab:integrity-parts}.

\begin{landscape}
\begin{figure}[p]
\centering
\IfFileExists{integrity_architecture.pdf}{%
  \includegraphics[width=\linewidth,height=0.86\textheight,keepaspectratio]{integrity_architecture.pdf}}{%
\IfFileExists{integrity_architecture.png}{%
  \includegraphics[width=\linewidth,height=0.86\textheight,keepaspectratio]{integrity_architecture.png}}{%
  \fbox{\parbox{0.8\linewidth}{\centering\vspace{2cm}
  Экспортируйте схему из draw.io в файл \texttt{integrity\_architecture.pdf}.\vspace{2cm}}}}}
\caption{Архитектура подсистемы контроля целостности}\label{fig:integrity-arch}
\end{figure}
\end{landscape}

{\small
\begin{longtable}{|L{3.4cm}|L{2.6cm}|L{8.2cm}|}
\caption{Компоненты подсистемы контроля целостности}\label{tab:integrity-parts}\\
\hline
\textbf{Компонент} & \textbf{Модуль} & \textbf{Назначение}\\
\hline\endfirsthead
\hline
\textbf{Компонент} & \textbf{Модуль} & \textbf{Назначение}\\
\hline\endhead
API-сервер & \texttt{protocol} &
  UNIX-сокет \texttt{/run/fw/integrity.sock}, JSON-строки (АПИ-03, АПИ-04),
  коды ошибок (АПИ-09). Права проверяются по контексту вызова: ТЧ, ЧЗ или
  подсистема; отказ регистрируется событием~106\\ \hline
Планировщик & \texttt{daemon} &
  Проверка при запуске до сигнала готовности systemd, каждые
  1\,--\,1440~мин (по умолчанию 10) и по запросу (ЦЛС-03)\\ \hline
Ядро & \texttt{service} &
  Команды API, режимы «обычный» и «аварийный», легитимные изменения
  файлов (ЦЛС-02), период проверки. Проверка, восстановление и обновление
  эталонов выполняются по очереди под одной блокировкой\\ \hline
Проверка файлов & \texttt{checker}, \texttt{manifest} &
  SHA-256 каждого файла с эталоном; поиск изменённых, удалённых,
  посторонних (в том числе скрытых) и нечитаемых файлов (ЦЛС-01)\\ \hline
Проверка ОС & \texttt{os\_state} &
  Сверка с таблицей ОБР-02: пересылка ядром, IPv6, адреса портов, promisc
  и MTU, мосты и VLAN, аппаратные функции, правило \texttt{drop} на
  \texttt{netdev ingress}\\ \hline
Восстановление & \texttt{store} &
  Изменённый или удалённый файл восстанавливается из резервной копии
  после её проверки по эталону; посторонний файл перемещается в карантин
  (ЦЛС-05)\\ \hline
Глобальная блокировка & \texttt{actions} &
  Команда \texttt{set\_emergency} подсистеме фильтрации и остановка
  служб fw-filter, fw-web, ssh (ЦЛС-04); при выходе~— обратные действия
  (ЦЛС-06)\\ \hline
Отправка событий & \texttt{events} &
  Команда \texttt{log\_event} подсистеме регистрации (коды 416,
  500\,--\,505); очередь до 1000 событий, если регистрация недоступна\\ \hline
\end{longtable}
}

\textbf{Режимы работы.} В обычном режиме трафик обрабатывается, а файлы
проверяются по расписанию. Любое нарушение (или невозможность проверки:
нет эталонов, повреждён манифест) переводит подсистему в аварийный режим.
Сначала режим записывается в \texttt{state.json}, затем включается
глобальная блокировка, затем регистрируются события~501 и~502.
Возврат в обычный режим происходит только после повторной проверки без
нарушений (событие~504). Запуск fw-filter и fw-web защищён проверкой
\texttt{fw-integrity gate} в \texttt{ExecStartPre}: служба стартует, только
если в текущей загрузке уже была успешная проверка (ЦЛС-07).

\textbf{Легитимные изменения.} Подсистемы, которые пишут конфигурационные
файлы (политику, сохранённую конфигурацию, учётные записи), вызывают
\texttt{begin\_update}, атомарно записывают файл и вызывают
\texttt{commit\_update}. Пока изменение не завершено, файл не проверяется.
После \texttt{commit\_update} эталон и резервная копия обновляются
(событие~505). Изменение без этих вызовов считается нарушением.

\textbf{Данные подсистемы} хранятся в каталоге
\texttt{/var/lib/fw/integrity}, доступном только \texttt{root}:
\begin{itemize}
\item \texttt{reference.json}~— эталонные SHA-256, вид, размер и права файлов;
\item \texttt{backup/}~— резервные копии, имя файла~— его SHA-256;
\item \texttt{quarantine/}~— посторонние файлы, убранные при восстановлении;
\item \texttt{state.json}~— режим, перечень нарушений, период и
  идентификатор загрузки.
\end{itemize}
Перечень контролируемых файлов задаёт манифест \texttt{manifest.json}:
\texttt{static} для кода и неизменяемых файлов, \texttt{config} для
изменяемых конфигураций. Сам манифест и файл настроек всегда
контролируются.

%  Рисунок: в draw.io открыть integrity_architecture.drawio,
%  «Файл -> Экспорт -> PDF» с галочкой «Обрезать», сохранить как
%  integrity_architecture.pdf рядом с документом.
% =====================================================================
\section{Архитектура подсистемы контроля целостности}\label{sec:integrity-arch}

Подсистема реализована процессом \texttt{fw-integrity} на Python~3.13,
используется только стандартная библиотека. Процесс работает от
\texttt{root}: ему нужно читать все контролируемые файлы, восстанавливать
их и останавливать службы в аварийном режиме. Схема приведена на
рисунке~\ref{fig:integrity-arch}, назначение компонентов~— в
таблице~\ref{tab:integrity-parts}.

\begin{landscape}
\begin{figure}[p]
\centering
\IfFileExists{integrity_architecture.pdf}{%
  \includegraphics[width=\linewidth,height=0.86\textheight,keepaspectratio]{integrity_architecture.pdf}}{%
\IfFileExists{integrity_architecture.png}{%
  \includegraphics[width=\linewidth,height=0.86\textheight,keepaspectratio]{integrity_architecture.png}}{%
  \fbox{\parbox{0.8\linewidth}{\centering\vspace{2cm}
  Экспортируйте схему из draw.io в файл \texttt{integrity\_architecture.pdf}.\vspace{2cm}}}}}
\caption{Архитектура подсистемы контроля целостности}\label{fig:integrity-arch}
\end{figure}
\end{landscape}

{\small
\begin{longtable}{|L{3.4cm}|L{2.6cm}|L{8.2cm}|}
\caption{Компоненты подсистемы контроля целостности}\label{tab:integrity-parts}\\
\hline
\textbf{Компонент} & \textbf{Модуль} & \textbf{Назначение}\\
\hline\endfirsthead
\hline
\textbf{Компонент} & \textbf{Модуль} & \textbf{Назначение}\\
\hline\endhead
API-сервер & \texttt{protocol} &
  UNIX-сокет \texttt{/run/fw/integrity.sock}, JSON-строки (АПИ-03, АПИ-04),
  коды ошибок (АПИ-09). Права проверяются по контексту вызова: ТЧ, ЧЗ или
  подсистема; отказ регистрируется событием~106\\ \hline
Планировщик & \texttt{daemon} &
  Проверка при запуске до сигнала готовности systemd, каждые
  1\,--\,1440~мин (по умолчанию 10) и по запросу (ЦЛС-03)\\ \hline
Ядро & \texttt{service} &
  Команды API, режимы «обычный» и «аварийный», легитимные изменения
  файлов (ЦЛС-02), период проверки. Проверка, восстановление и обновление
  эталонов выполняются по очереди под одной блокировкой\\ \hline
Проверка файлов & \texttt{checker}, \texttt{manifest} &
  SHA-256 каждого файла с эталоном; поиск изменённых, удалённых,
  посторонних (в том числе скрытых) и нечитаемых файлов (ЦЛС-01)\\ \hline
Проверка ОС & \texttt{os\_state} &
  Сверка с таблицей ОБР-02: пересылка ядром, IPv6, адреса портов, promisc
  и MTU, мосты и VLAN, аппаратные функции, правило \texttt{drop} на
  \texttt{netdev ingress}\\ \hline
Восстановление & \texttt{store} &
  Изменённый или удалённый файл восстанавливается из резервной копии
  после её проверки по эталону; посторонний файл перемещается в карантин
  (ЦЛС-05)\\ \hline
Глобальная блокировка & \texttt{actions} &
  Команда \texttt{set\_emergency} подсистеме фильтрации и остановка
  служб fw-filter, fw-web, ssh (ЦЛС-04); при выходе~— обратные действия
  (ЦЛС-06)\\ \hline
Отправка событий & \texttt{events} &
  Команда \texttt{log\_event} подсистеме регистрации (коды 416,
  500\,--\,505); очередь до 1000 событий, если регистрация недоступна\\ \hline
\end{longtable}
}

\textbf{Режимы работы.} В обычном режиме трафик обрабатывается, а файлы
проверяются по расписанию. Любое нарушение (или невозможность проверки:
нет эталонов, повреждён манифест) переводит подсистему в аварийный режим.
Сначала режим записывается в \texttt{state.json}, затем включается
глобальная блокировка, затем регистрируются события~501 и~502.
Возврат в обычный режим происходит только после повторной проверки без
нарушений (событие~504). Запуск fw-filter и fw-web защищён проверкой
\texttt{fw-integrity gate} в \texttt{ExecStartPre}: служба стартует, только
если в текущей загрузке уже была успешная проверка (ЦЛС-07).

\textbf{Легитимные изменения.} Подсистемы, которые пишут конфигурационные
файлы (политику, сохранённую конфигурацию, учётные записи), вызывают
\texttt{begin\_update}, атомарно записывают файл и вызывают
\texttt{commit\_update}. Пока изменение не завершено, файл не проверяется.
После \texttt{commit\_update} эталон и резервная копия обновляются
(событие~505). Изменение без этих вызовов считается нарушением.

\textbf{Данные подсистемы} хранятся в каталоге
\texttt{/var/lib/fw/integrity}, доступном только \texttt{root}:
\begin{itemize}
\item \texttt{reference.json}~— эталонные SHA-256, вид, размер и права файлов;
\item \texttt{backup/}~— резервные копии, имя файла~— его SHA-256;
\item \texttt{quarantine/}~— посторонние файлы, убранные при восстановлении;
\item \texttt{state.json}~— режим, перечень нарушений, период и
  идентификатор загрузки.
\end{itemize}
Перечень контролируемых файлов задаёт манифест \texttt{manifest.json}:
\texttt{static} для кода и неизменяемых файлов, \texttt{config} для
изменяемых конфигураций. Сам манифест и файл настроек всегда
контролируются.

%  Рисунок: в draw.io открыть integrity_architecture.drawio,
%  «Файл -> Экспорт -> PDF» с галочкой «Обрезать», сохранить как
%  integrity_architecture.pdf рядом с документом.
% =====================================================================
\section{Архитектура подсистемы контроля целостности}\label{sec:integrity-arch}

Подсистема реализована процессом \texttt{fw-integrity} на Python~3.13,
используется только стандартная библиотека. Процесс работает от
\texttt{root}: ему нужно читать все контролируемые файлы, восстанавливать
их и останавливать службы в аварийном режиме. Схема приведена на
рисунке~\ref{fig:integrity-arch}, назначение компонентов~— в
таблице~\ref{tab:integrity-parts}.

\begin{landscape}
\begin{figure}[p]
\centering
\IfFileExists{integrity_architecture.pdf}{%
  \includegraphics[width=\linewidth,height=0.86\textheight,keepaspectratio]{integrity_architecture.pdf}}{%
\IfFileExists{integrity_architecture.png}{%
  \includegraphics[width=\linewidth,height=0.86\textheight,keepaspectratio]{integrity_architecture.png}}{%
  \fbox{\parbox{0.8\linewidth}{\centering\vspace{2cm}
  Экспортируйте схему из draw.io в файл \texttt{integrity\_architecture.pdf}.\vspace{2cm}}}}}
\caption{Архитектура подсистемы контроля целостности}\label{fig:integrity-arch}
\end{figure}
\end{landscape}

{\small
\begin{longtable}{|L{3.4cm}|L{2.6cm}|L{8.2cm}|}
\caption{Компоненты подсистемы контроля целостности}\label{tab:integrity-parts}\\
\hline
\textbf{Компонент} & \textbf{Модуль} & \textbf{Назначение}\\
\hline\endfirsthead
\hline
\textbf{Компонент} & \textbf{Модуль} & \textbf{Назначение}\\
\hline\endhead
API-сервер & \texttt{protocol} &
  UNIX-сокет \texttt{/run/fw/integrity.sock}, JSON-строки (АПИ-03, АПИ-04),
  коды ошибок (АПИ-09). Права проверяются по контексту вызова: ТЧ, ЧЗ или
  подсистема; отказ регистрируется событием~106\\ \hline
Планировщик & \texttt{daemon} &
  Проверка при запуске до сигнала готовности systemd, каждые
  1\,--\,1440~мин (по умолчанию 10) и по запросу (ЦЛС-03)\\ \hline
Ядро & \texttt{service} &
  Команды API, режимы «обычный» и «аварийный», легитимные изменения
  файлов (ЦЛС-02), период проверки. Проверка, восстановление и обновление
  эталонов выполняются по очереди под одной блокировкой\\ \hline
Проверка файлов & \texttt{checker}, \texttt{manifest} &
  SHA-256 каждого файла с эталоном; поиск изменённых, удалённых,
  посторонних (в том числе скрытых) и нечитаемых файлов (ЦЛС-01)\\ \hline
Проверка ОС & \texttt{os\_state} &
  Сверка с таблицей ОБР-02: пересылка ядром, IPv6, адреса портов, promisc
  и MTU, мосты и VLAN, аппаратные функции, правило \texttt{drop} на
  \texttt{netdev ingress}\\ \hline
Восстановление & \texttt{store} &
  Изменённый или удалённый файл восстанавливается из резервной копии
  после её проверки по эталону; посторонний файл перемещается в карантин
  (ЦЛС-05)\\ \hline
Глобальная блокировка & \texttt{actions} &
  Команда \texttt{set\_emergency} подсистеме фильтрации и остановка
  служб fw-filter, fw-web, ssh (ЦЛС-04); при выходе~— обратные действия
  (ЦЛС-06)\\ \hline
Отправка событий & \texttt{events} &
  Команда \texttt{log\_event} подсистеме регистрации (коды 416,
  500\,--\,505); очередь до 1000 событий, если регистрация недоступна\\ \hline
\end{longtable}
}

\textbf{Режимы работы.} В обычном режиме трафик обрабатывается, а файлы
проверяются по расписанию. Любое нарушение (или невозможность проверки:
нет эталонов, повреждён манифест) переводит подсистему в аварийный режим.
Сначала режим записывается в \texttt{state.json}, затем включается
глобальная блокировка, затем регистрируются события~501 и~502.
Возврат в обычный режим происходит только после повторной проверки без
нарушений (событие~504). Запуск fw-filter и fw-web защищён проверкой
\texttt{fw-integrity gate} в \texttt{ExecStartPre}: служба стартует, только
если в текущей загрузке уже была успешная проверка (ЦЛС-07).

\textbf{Легитимные изменения.} Подсистемы, которые пишут конфигурационные
файлы (политику, сохранённую конфигурацию, учётные записи), вызывают
\texttt{begin\_update}, атомарно записывают файл и вызывают
\texttt{commit\_update}. Пока изменение не завершено, файл не проверяется.
После \texttt{commit\_update} эталон и резервная копия обновляются
(событие~505). Изменение без этих вызовов считается нарушением.

\textbf{Данные подсистемы} хранятся в каталоге
\texttt{/var/lib/fw/integrity}, доступном только \texttt{root}:
\begin{itemize}
\item \texttt{reference.json}~— эталонные SHA-256, вид, размер и права файлов;
\item \texttt{backup/}~— резервные копии, имя файла~— его SHA-256;
\item \texttt{quarantine/}~— посторонние файлы, убранные при восстановлении;
\item \texttt{state.json}~— режим, перечень нарушений, период и
  идентификатор загрузки.
\end{itemize}
Перечень контролируемых файлов задаёт манифест \texttt{manifest.json}:
\texttt{static} для кода и неизменяемых файлов, \texttt{config} для
изменяемых конфигураций. Сам манифест и файл настроек всегда
контролируются.
